\documentclass[10pt,a4paper]{amsart}
\usepackage{amsmath,amssymb,amsthm,graphicx}
% Retained minimal preamble for the Stacks Project chapter more-morphisms.tex (More on Morphisms).
% Verbatim excerpt of the chapter preamble: Xy-pic setup (native preamble lines 14-18) and the native
% theorem declarations with the mathematical macros (native preamble lines 154-217). Copied mechanically;
% nothing added, removed or redefined. The wrapper class and the amsmath/amssymb/amsthm/graphicx base
% packages are supplied by the assembler, so the class/externaldocument part of the native preamble is omitted.
% For commutative diagrams we use Xy-pic
\usepackage[all]{xy}

% We use 2cell for 2-commutative diagrams.
\xyoption{2cell}
% Theorem environments.
%
\theoremstyle{plain}
\newtheorem{theorem}[subsection]{Theorem}
\newtheorem{proposition}[subsection]{Proposition}
\newtheorem{lemma}[subsection]{Lemma}

\theoremstyle{definition}
\newtheorem{definition}[subsection]{Definition}
\newtheorem{example}[subsection]{Example}
\newtheorem{exercise}[subsection]{Exercise}
\newtheorem{situation}[subsection]{Situation}

\theoremstyle{remark}
\newtheorem{remark}[subsection]{Remark}
\newtheorem{remarks}[subsection]{Remarks}

\numberwithin{equation}{subsection}

% Macros
%
\def\lim{\mathop{\mathrm{lim}}\nolimits}
\def\colim{\mathop{\mathrm{colim}}\nolimits}
\def\Spec{\mathop{\mathrm{Spec}}}
\def\Hom{\mathop{\mathrm{Hom}}\nolimits}
\def\Ext{\mathop{\mathrm{Ext}}\nolimits}
\def\SheafHom{\mathop{\mathcal{H}\!\mathit{om}}\nolimits}
\def\SheafExt{\mathop{\mathcal{E}\!\mathit{xt}}\nolimits}
\def\Sch{\mathit{Sch}}
\def\Mor{\mathop{\mathrm{Mor}}\nolimits}
\def\Ob{\mathop{\mathrm{Ob}}\nolimits}
\def\Sh{\mathop{\mathit{Sh}}\nolimits}
\def\NL{\mathop{N\!L}\nolimits}
\def\CH{\mathop{\mathrm{CH}}\nolimits}
\def\proetale{{pro\text{-}\acute{e}tale}}
\def\etale{{\acute{e}tale}}
\def\QCoh{\mathit{QCoh}}
\def\Ker{\mathop{\mathrm{Ker}}}
\def\Im{\mathop{\mathrm{Im}}}
\def\Coker{\mathop{\mathrm{Coker}}}
\def\Coim{\mathop{\mathrm{Coim}}}

% Boxtimes
%
\DeclareMathSymbol{\boxtimes}{\mathbin}{AMSa}{"02}

%
% Macros for moduli stacks/spaces
%
\def\QCohstack{\mathcal{QC}\!\mathit{oh}}
\def\Cohstack{\mathcal{C}\!\mathit{oh}}
\def\Spacesstack{\mathcal{S}\!\mathit{paces}}
\def\Quotfunctor{\mathrm{Quot}}
\def\Hilbfunctor{\mathrm{Hilb}}
\def\Curvesstack{\mathcal{C}\!\mathit{urves}}
\def\Polarizedstack{\mathcal{P}\!\mathit{olarized}}
\def\Complexesstack{\mathcal{C}\!\mathit{omplexes}}
% \Pic is the operator that assigns to X its picard group, usage \Pic(X)
% \Picardstack_{X/B} denotes the Picard stack of X over B
% \Picardfunctor_{X/B} denotes the Picard functor of X over B
\def\Pic{\mathop{\mathrm{Pic}}\nolimits}
\def\Picardstack{\mathcal{P}\!\mathit{ic}}
\def\Picardfunctor{\mathrm{Pic}}
\def\Deformationcategory{\mathcal{D}\!\mathit{ef}}
\usepackage{cleveref}
\crefname{originalreference}{original reference}{original references}
\Crefname{originalreference}{Original reference}{Original references}
\makeatletter
\newlabel{descent-lemma-descending-property-finite}{{02LA}{source}{}{sourcepointer-1}{}}
\newlabel{descent-lemma-descending-property-finite@cref}{{[originalreference][0][]02LA}{[1][0][]source}{}{}{}}
\makeatother
\crefname{originallemma}{original lemma}{original lemmas}
\Crefname{originallemma}{Original lemma}{Original lemmas}
\makeatletter
\newlabel{lemma-etale-splits-off-quasi-finite-part-technical-variant}{{02LO}{source}{}{sourcepointer-2}{}}
\newlabel{lemma-etale-splits-off-quasi-finite-part-technical-variant@cref}{{[originallemma][0][]02LO}{[1][0][]source}{}{}{}}
\makeatother
\crefname{originallemma}{original lemma}{original lemmas}
\Crefname{originallemma}{Original lemma}{Original lemmas}
\makeatletter
\newlabel{lemma-semicontinuous-int-w}{{0F3J}{source}{}{sourcepointer-3}{}}
\newlabel{lemma-semicontinuous-int-w@cref}{{[originallemma][0][]0F3J}{[1][0][]source}{}{}{}}
\makeatother
\title{Where a positive weighting attains its maximum the morphism is finite}
\author{The Stacks Project Authors}
\date{stacks-project:more-morphisms.tex:local-snapshot-20261011; GNU FDL 1.2 or later}
\begin{document}
\maketitle
\noindent\textit{Editorial note.} Counted span is the author's own proof environment, source lines 22757-22788, immediately following the statement at 22748-22755.\par\medskip
\noindent\textit{Editorial note.} One cross-chapter pointer (Descent) plus two same-chapter lemmas are recorded in assembly.external\_references.\par\medskip
\noindent\textit{Editorial note.} The displayed decomposition of the integral is retained verbatim; the conclusion W empty is drawn inside the span.\par\medskip
\noindent\textit{Allowed standard prerequisites (editorial).}\begin{itemize}
\item Upper semi-continuity of the integral. chapter lemma lemma-semicontinuous-int-w, cited at proof line 22758. Shows the locus where the integral equals its maximum is open.
\item Etale splitting off of the quasi-finite part. chapter lemma lemma-etale-splits-off-quasi-finite-part-technical-variant, cited at proof lines 22760-22761. Provides the etale neighbourhood and the decomposition used to write the integral as a sum of a finite piece and the value at the point.
\item Finiteness descends along the open etale morphism. Descent, lemma descent-lemma-descending-property-finite, cited at proof line 22783. Transfers finiteness from the etale neighbourhood back to f over the image of the open morphism.
\end{itemize}\medskip
\ifdefined\hypertarget\hypertarget{sourcepointer-1}{}\fi
\noindent\textit{Source reference.} Original reference 02LA: focused-seeds/geometry/S9-morph-max-int-finite/ref-tags.txt, source lines 4--4. Navigation pointer only: the counted proof cites this statement by \textbackslash{}ref as a standing input; the deduction the counted proof has to make is carried out inside the counted span, and no part of that deduction is delegated to this pointer.\par\medskip
\ifdefined\hypertarget\hypertarget{sourcepointer-2}{}\fi
\noindent\textit{Source reference.} Original lemma 02LO: sources/stacks-more-morphisms/pinned-original-source.tex, source lines 12136--12136. Navigation pointer only: the counted proof cites this statement by \textbackslash{}ref as a standing input; the deduction the counted proof has to make is carried out inside the counted span, and no part of that deduction is delegated to this pointer.\par\medskip
\ifdefined\hypertarget\hypertarget{sourcepointer-3}{}\fi
\noindent\textit{Source reference.} Original lemma 0F3J: sources/stacks-more-morphisms/pinned-original-source.tex, source lines 22695--22695. Navigation pointer only: the counted proof cites this statement by \textbackslash{}ref as a standing input; the deduction the counted proof has to make is carried out inside the counted span, and no part of that deduction is delegated to this pointer.\par\medskip
\begin{lemma}
\label{lemma-max-int-finite}
Let $f : X \to Y$ be a separated, locally quasi-finite morphism.
Let $w : X \to \mathbf{Z}_{> 0}$ be a positive weighting of $f$.
Assume $\int_w f$ attains its maximum $d$ and let $Y_d \subset Y$
be the open set of points $y$ with $(\int_f w)(y) = d$. Then
the morphism $f^{-1}(Y_d) \to Y_d$ is finite.
\end{lemma}

\begin{proof}
Observe that $Y_d$ is open by Lemma \ref{lemma-semicontinuous-int-w}.
Let $y \in Y_d$. Say $x_1, \ldots, x_n$ are the points of $X$
lying over $y$. Apply
Lemma \ref{lemma-etale-splits-off-quasi-finite-part-technical-variant}
to get an \'etale neighbourhood $(U, u) \to (Y, y)$ and a decomposition
$$
U \times_Y X =
W \amalg
\ \coprod\nolimits_{i = 1, \ldots, n}
\ \coprod\nolimits_{j = 1, \ldots, m_i}
V_{i, j}
$$
as in locus citatus. Observe that $d = \sum w(v_{i, j})$ where
$w(v_{i, j}) = w(x_i)$. Since $\int_{V_{i, j} \to U} w|_{V_{i, j}}$
is locally constant by definition, we may after shrinking $U$
assume these functions are constant with value $w(v_{i, j})$.
We conclude that
$$
\textstyle{\int}_{U \times_Y X \to U} w|_{U \times_Y X} =
\textstyle{\int}_{W \to U} w|_W +
\sum \textstyle{\int}_{V_{i, j} \to U} w|_{V_{i, j}} =
\textstyle{\int}_{W \to U} w|_W + (\int_f w)(y)
$$
This is $\geq (\int_f w)(y) = d$ and we conclude that $W$
must be the emptyset. Thus $U \times_Y X \to U$ is finite.
By Descent, Lemma \ref{descent-lemma-descending-property-finite}
this implies that $X \to Y$ is finite over
the image of the open morphism $U \to Y$. In other words,
we see that $f$ is finite over an open neighbourhood of $y$
as desired.
\end{proof}

\end{document}
